\begin{center}
\begin{circuitikz}[european resistors]
 \draw (0,0) to[sV, l={$V$}] (0,2.5)
   to[short, i={$I_\mathrm{I}$}] (1,2.5)
   to[R, l={$R_1$}] (3,2.5)
   to[R, l={$R_2$}] (5,2.5) -- (5,0) -- (0,0);
 \draw (7,0) to[sV, l={$V$}] (7,2.5)
   to[short, i={$I_\mathrm{II}$}] (8,2.5) -- (11,2.5);
 \draw (7,0) -- (11,0);
 \draw (9,2.5) to[R, l={$R_1$}, *-*] (9,0);
 \draw (11,2.5) to[R, l={$R_2$}] (11,0);
 \node at (2.5,-0.6) {setting I};
 \node at (9,-0.6) {setting II};
\end{circuitikz}
\end{center}

Setting I, in series: the resistances add up.
\begin{align}
 I_\mathrm{I} &= \IsF \\
   &= \frac{\V}{\Ra+\Rb} \\
   &= \Is \approx \result{\IsP{0}{2}}
\end{align}

Setting II, in parallel: each wire is at the full mains voltage, and the currents add up.
\begin{align}
 I_\mathrm{II} &= \IpF \\
   &= \frac{\V}{\Ra}+\frac{\V}{\Rb} \\
   &= \result{\IpP{0}{2}}
\end{align}
In parallel, the current is $\ratioP{0}{2}$ times larger: setting II heats much more.
